\subsection{Actual dispatch and timing-aware attacks}\label{sec:cover-study}
\textbf{Immediate controller cover is removable because it can skip the model wait preceding honest calls.}
Two fixed-weight planners, Qwen3:8B and Llama3.1:8B, each execute 300 fixtures in five conditions, yielding 3,000 main episodes and 600 additional shadow chains. The six business actions use seven MCP services through LangGraph over TLS~1.3/HTTP/2. Original partitions contain 120 training, 60 validation and 120 paired test fixtures per cell. The task label is requested intent; failures remain included. Cover is isolated from real business state and hidden from the real planner.

No cover establishes the baseline. Immediate cover dispatches after an honest response. Historical sequential and concurrent policies sample training-only honest gaps and service durations, wait in real time and hold responses for the sampled minimum duration. The concurrent dispatcher overlaps real model inference. Shadow runs a complete second planner on an independent six-class ticket; either chain can start first. Calls to the same service serialize. The original observer receives destinations, ciphertext bytes, TLS record completion times and supplied episode/call windows. Section~\ref{sec:no-call} removes the last of these call-level privileges.

\begin{table*}[t]\centering\small
\caption{Joint evaluation of the five live policies, 120 test fixtures per row; uniform six-class reference 20/120. Sel. is the validation-selected non-timing/adaptive count; Max is the largest test count in the frozen suite, reported descriptively, not substituted into the primary tests. GMM is role-label-free removal. Bytes gives whole-connection/call-window factors; Tokens includes both planners. Goal and State are separate counts. In shadow, both selectors choose the identical non-timing attack.}\label{tab:live-timing-main}
\setlength{\tabcolsep}{4pt}
\begin{tabularx}{\textwidth}{@{}lYrrrrrr@{}}\toprule
Planner & Policy & Sel. S/A & Max & GMM & Bytes$\times$ & Tokens$\times$ & Goal/State \\\midrule

Qwen3:8B & No cover & 117/117 & 117 & -- & 1.00/1.00 & 1.00 & 113/114 \\
Qwen3:8B & Immediate & 75/112 & 112 & 112 & 1.15/1.33 & 1.00 & 111/112 \\
Qwen3:8B & History, sequential & 76/93 & 99 & 53 & 1.15/1.32 & 1.00 & 112/113 \\
Qwen3:8B & History, concurrent & 66/99 & 99 & 38 & 1.15/1.33 & 1.00 & 111/112 \\
Qwen3:8B & Planner shadow & 54/54 & 70 & 34 & 1.31/1.68 & 2.01 & 112/113 \\
Llama3.1:8B & No cover & 110/110 & 112 & -- & 1.00/1.00 & 1.00 & 82/70 \\
Llama3.1:8B & Immediate & 80/108 & 108 & 108 & 1.13/1.27 & 1.01 & 83/67 \\
Llama3.1:8B & History, sequential & 82/104 & 104 & 14 & 1.13/1.26 & 1.01 & 84/67 \\
Llama3.1:8B & History, concurrent & 68/106 & 106 & 66 & 1.13/1.27 & 1.02 & 82/67 \\
Llama3.1:8B & Planner shadow & 64/64 & 64 & 43 & 1.35/1.71 & 1.99 & 82/67 \\
\bottomrule\end{tabularx}\end{table*}


\paragraph{Removal and adaptation}
An unlabelled training-gap mixture selects the short-gap component for removal; a no-cover-trained task classifier decodes the remaining trace. In the immediate arm this restores 112/120 and 108/120 correct predictions, versus selected non-timing references 75 and 80. Both meet the prespecified replication criterion and pass Holm correction. Independent TCP reassembly matches all 773,450 recorded TLS record types, lengths and counts. Recomputing the GMM attack with packet-capture clocks reproduces the scores; software completion timestamps retain measured error (Supplement S2).

Historical sampling changes the short-gap attack, but adapted context removal or joint decoding still recognizes 93/120 and 104/120 sequential executions. The Qwen comparison with its 76/120 non-timing reference has Holm $p=.0559$; Llama's comparison with 82/120 has $p=.00078$. Concurrent adapted counts are 99/120 and 106/120, against 66 and 68. Their 27.5- and 31.7-point \emph{attacker advantages under the same policy} both pass correction. The suite maxima in Table~\ref{tab:live-timing-main} additionally expose Qwen's sequential joint decoder at 99/120, rather than replace the frozen selection after testing.

The policy samples minimum waits, not achieved timing distributions. Qwen's mean 40.5\,ms sampled hold becomes 100.4\,ms sequential and 171.2\,ms concurrent at the record observer. A failed simple mixture is therefore not evidence of successful concealment. In the saved packet-clock diagnostic, Llama's sequential mixture deletes 477/634 honest calls while detecting only 64/182 cover calls, explaining its 14/120 task count. Adapted joint decoding on the same condition reaches 104/120.

\paragraph{Planner-generated cover and cost}
In the call-window suite, shadow's strongest attack is the size/destination RF: 70/120 Qwen and 64/120 Llama, or 58.3\% and 53.3\%, versus the 16.7\% uniform reference. A chain-separation attack scores 49 and 34 and exactly recovers the real chain in 29 and 32 fixtures. The validation selectors choose the same non-timing attack for both roles; their equality is not an independent role-indistinguishability test. Shadow uses 2.01/1.99 times model tokens, 1.68/1.71 times call-window bytes and mean real-chain times 4.80/3.94\,s versus baseline 2.91/3.33\,s. It is a measured operating point, not an equal-budget demonstration that planning itself is superior. Supplement S2 retains every attack, primary test and cost.
